\chapter{Riferimento rapido}
\label{ch:ref}

\section{Opcode SPI}
\begin{tabularx}{\textwidth}{C{1.4cm} L{3.2cm} C{2.4cm} Y}
\toprule
\rowh \thd{Valore} & \thd{Nome} & \thd{Risposta} & \thd{Sintesi} \\
\midrule
\op{0x00} & NOP & --- & idle \\
\rowa \op{0x01} & WRITE\_RAM & --- & scrittura blocco PSRAM \\
\op{0x02} & READ\_RAM & \code{len} B & lettura blocco PSRAM \\
\rowa \op{0x0F} & RESET & --- & reset motore + latch STATUS \\
\op{0x10} & SET\_BASE & --- & imposta base/registro (sel 0..10) \\
\rowa \op{0x11} & SET\_NET\_TYPE & --- & tipo rete \#1/\#2 \\
\rowa \op{0x20} & START & --- & avvio single-layer \\
\op{0x21} & STATUS & 1 B & busy(live)/done(sticky) \\
\rowa \op{0x22} & READ\_OUTPUT & N\_NEURONS B & \code{y\_bus} \\
\op{0x23} & RUN\_NETWORK & --- & avvio multi-layer \\
\rowa \op{0x30} & READ\_CONFIG & 11 B & record configurazione \\
\bottomrule
\end{tabularx}

\section{Byte di STATUS}
\begin{center}
\begin{tikzpicture}[font=\scriptsize]
\foreach \i/\lbl [count=\x from 0] in {7/0,6/0,5/0,4/0,3/0,2/0,1/{done},0/{busy}}{
  \node[fnreg,minimum width=13mm,minimum height=9mm] (b\x) at (\x*13mm,0) {\lbl};
  \node[font=\tiny,text=fnGrey,above=0.5mm of b\x] {bit \i};
}
\node[fill=fnAmber,text=white,rounded corners=1pt,inner sep=1.5pt,font=\tiny]
   at (b7.center){riservati = 0};
\node[fill=fnTeal,text=white,rounded corners=1pt,inner sep=1.5pt,font=\tiny]
   at (b6.center){};
\end{tikzpicture}
\end{center}
\code{done} è sticky, clear-on-read; \code{busy} è live; \code{bit2=err} (guard grafo).

\section{Selettori SET\_BASE}
\begin{multicols}{2}\footnotesize
\begin{itemize}
\item 0 --- \code{x\_base}
\item 1 --- \code{w\_base}
\item 2 --- \code{bias\_addr}
\item 3 --- \code{table\_base}
\item 4 --- \code{buf\_a\_base}
\columnbreak
\item 5 --- \code{buf\_b\_base}
\item 6 --- \code{activation} (single-layer)
\item 7 --- \code{n\_inputs\_real} (single-layer)
\item 8 --- \code{n\_neurons\_real} (single-layer)
\item 9 --- \code{num\_neurons\_graph} (Tipo \#2)
\item 10 --- \code{n\_out} (Tipo \#2)
\end{itemize}
\end{multicols}

\section{Tabella descrittori (11 byte/layer, MSB-first)}
\begin{center}
\begin{tikzpicture}[font=\scriptsize,node distance=0mm]
 \node[fnreg,minimum width=20mm,minimum height=8mm](a){\code{w\_base}\\3B};
 \node[fnreg,minimum width=20mm,minimum height=8mm,right=0mm of a](b){\code{bias\_addr}\\3B};
 \node[fnreg,minimum width=14mm,minimum height=8mm,right=0mm of b](c){\code{act}\\1B};
 \node[fnreg,minimum width=22mm,minimum height=8mm,right=0mm of c](d){\code{n\_inputs\_real}\\2B};
 \node[fnreg,minimum width=22mm,minimum height=8mm,right=0mm of d](e){\code{n\_neurons\_real}\\2B};
\end{tikzpicture}
\end{center}

\section{Parametri di build}
\begin{multicols}{2}\footnotesize
\begin{itemize}
\item \code{DATA\_WIDTH} --- 8 (INT8)
\item \code{ACC\_WIDTH} --- 32 (INT32)
\item \code{N\_INPUTS} --- max ingressi
\item \code{N\_NEURONS} --- max neuroni
\item \code{PARALLEL} --- MAC simultanei
\columnbreak
\item \code{N\_LAYERS} --- max layer
\item \code{ADDR\_WIDTH} --- 23 (8 MB)
\item \code{MEM\_DATA\_WIDTH} --- 16
\item \code{CLK\_FREQ\_MHZ} --- timing PSRAM
\end{itemize}
\end{multicols}
